% !TeX root = ../collection-solutions.tex

\titledquestion{Why are the wealthiest so wealthy?}

\citet{hubmer2026wealthiest} use Norwegian administrative panel data to trace how the top 0.1\% accumulated their wealth over 26 years.

\begin{parts}
    \part[4] What is the key takeaway of the paper? Explain how the authors arrive at this conclusion: what is the decomposition method, what are the headline magnitudes, and how do the authors use these facts to discriminate between theories of wealth inequality?

    \begin{solutionorbox}[12cm]
        \emph{Takeaway} (2 points): The excess wealth of the top 0.1\% over the median household is accounted for by higher \emph{saving rates} ($\approx 36\%$), \emph{inheritances} ($\approx 31\%$), higher \emph{returns} ($\approx 28\%$) and higher labor income ($\approx 5\%$). These \emph{dynamic} facts reject all standard single-channel models of wealth inequality — even though each of them (suitably augmented) can match the cross-sectional concentration of wealth. Only a model combining entrepreneurship with decreasing returns to scale and non-homothetic (``capitalist spirit'') saving preferences fits the dynamics.

        \emph{How they get there} (2 points):
        \begin{itemize}
            \item Method: track households backward over 26 years through their \emph{budget identity} (wealth growth $=$ returns $+$ labor income $+$ inheritances, scaled by the saving rate). Replace one component at a time with the value of the middle class and recompute the counterfactual wealth path; average over all replacement orderings (Shapley--Owen) so the contributions sum to 100\%. This is an \emph{accounting} decomposition — it deliberately holds behavior fixed.
            \item A quarter of the top 0.1\% are ``New Money'': they start with roughly zero (on average negative) wealth and rise on extremely high returns and saving rates; ``Old Money'' heirs grow already-large fortunes mainly through inheritances and persistently high saving, with only modest returns.
            \item Model horse race: standard models (superstar earnings, heterogeneous returns, discount-factor heterogeneity, bequest motives) are each recalibrated to match cross-sectional wealth concentration, then confronted with the decomposition. Each fails in a characteristic way — e.g.\ the bequest and return-heterogeneity models attribute far too much to inheritances, the superstar model far too much to labor income and produces no Old Money.
        \end{itemize}
        Grading: 1 point for the four shares (approximate), 1 point for ``dynamics reject models that match the cross-section''; 1 point for the budget-identity/counterfactual-replacement method (incl.\ its accounting nature), 1 point for New vs.\ Old Money or a concrete model failure.
    \end{solutionorbox}

    \part[4] The data demand two things at once: self-made households that reach the top 0.1\% within $\sim$25 years (``New Money''), and heirs who keep growing large fortunes at high saving rates but moderate returns (``Old Money''). Explain intuitively how (i) entrepreneurship with decreasing returns to scale and a collateral constraint and (ii) non-homothetic ``capitalist spirit'' preferences jointly deliver both.

    \begin{solutionorbox}[12cm]
        \begin{itemize}
            \item (2 points) \emph{Decreasing returns + collateral constraint make returns non-monotone in wealth.} Across households, productivity and wealth are positively correlated, so \emph{average} returns rise with wealth (as in the data). But for a \emph{given} entrepreneur, the marginal return falls as the firm grows. A young, highly productive but collateral-constrained entrepreneur earns very high returns on a small base — the New Money rocket — while the same decreasing returns cap returns once the firm matures, so fortunes do not diverge without bound and established wealth earns only moderate returns (as Old Money does).
            \item (1 point) \emph{Non-homothetic preferences keep rich households saving.} With wealth in the utility function, the saving rate \emph{rises} with resources, so rich households — including heirs whose returns are unspectacular — keep saving large fractions of their income throughout life. This is what generates Old Money's high saving contribution, which returns alone cannot explain.
            \item (1 point) \emph{Timing matters:} unlike a warm-glow bequest motive, which only becomes relevant near the end of life (it is weighted by the probability of dying), the capitalist-spirit motive operates from the start — so \emph{young} wealthy households already save at high rates, as in the data. In a pure bequest model, young heirs would instead run down their wealth first.
        \end{itemize}
    \end{solutionorbox}
\end{parts}
